\section{Measurement and Evaluation Design}
\label{sec:design}
\subsection{Graph Construction and Views}
Graph construction defines which pages and links enter each measurement.
A node is a retained documentation page, and an arc is a resolved in-text link.
Duplicate arcs and self-loops are excluded. The binary adjacency matrix $A$ has
$A_{ij}=1$ when page $i$ links to page $j$. Graphs retain their frozen importer
and page inventories, including isolates; graph identity includes both nodes and
arcs. The corpus panel uses fourteen recovered non-FAQ snapshots.
Its extraction surface excludes website
navigation that the importer does not resolve. In Docker, for example, Hugo menus
and include shortcodes are outside the Markdown-link graph.

Measurements cover the full graph and its largest weakly connected component
(LWCC), where pages connect when link direction is ignored. Strong connectivity requires
directed paths in both directions. Ties in LWCC size are resolved by smallest node
identifier. The measured component's node count is reported beside the full
inventory, because selecting a component changes the object being described.
Each refined graph and null uses its own LWCC. The area comparison first induces
a graph on a common area node set and then selects that graph's LWCC; a common
input set does not imply identical measured nodes.

\subsection{Normalized Magnetic Profile}
\label{sec:magnetic-profile}
The magnetic profile retains link direction alongside structural comparators.
Let $W$ encode undirected interaction strength and $d_i$ its weighted degree:
\begin{equation}
 W=\tfrac12(A+A^{\mathsf T}),\qquad d_i=\sum_j W_{ij}.
\end{equation}
A one-way pair has weight $1/2$ and a reciprocal pair weight $1$.
Charge $q$ determines the rotation associated with orientation. With
$\theta_{ij}=2\pi q(A_{ij}-A_{ji})$ and
$H_{q,ij}=W_{ij}e^{i\theta_{ij}}$, define
\begin{equation}
 \mathcal L_q=\operatorname{diag}(\mathbf1_{d_i>0})-R H_q R,
 \quad R_{ii}=\begin{cases}d_i^{-1/2},&d_i>0,\\0,&d_i=0.\end{cases}
 \label{eq:operator}
\end{equation}
An isolate has a zero row and column, hence a zero eigenvalue. The diagonal is
not an unconditional identity. Eigenvalues are ordered increasingly,
$\lambda_1(q)\leq\lambda_2(q)\leq\cdots$; the profile retains the first magnetic
mode, which can be positive. At $q=0$, Eq.~\ref{eq:operator} uses the same $W$ and
becomes the weighted ordinary normalized Laplacian. On a connected view its
second eigenvalue measures weighted undirected connectivity.

The profile reports $\lambda_1$ and $\lambda_2$ at $q=.25$, with the grid
$q\in\{0,.05,.10,.15,.20,.25\}$ for sensitivity. This setting is held fixed
across comparisons. The structural comparators are matched $q=0$ connectivity, reciprocity,
out-degree Gini, size, fragmentation and generalized trophic incoherence.
The last quantity fits levels $h$ to unit increases along directed arcs:
\begin{equation}
 F_0=\min_h\frac{\sum_{ij}A_{ij}(h_j-h_i-1)^2}{\sum_{ij}A_{ij}}.
 \label{eq:trophic}
\end{equation}
Edgeless components do not enter this objective~\cite{mackay2020}.
For the separately identified historical refinement
test, the unnormalized operator is $L_q=D-H_q$, where $D=\operatorname{diag}(d)$.

\subsection{Ground-Mode Interpretation}
\label{sec:ground-mode}
Page mass and edge energy locate ground modes across the graph.
Let $U$ have orthonormal columns ($U^*U=I_k$) spanning the complete cluster of lowest-eigenvalue (ground) modes, of rank $k$.
Define page mass and edge energy by
\begin{align}
 m_i&=k^{-1}\sum_{r=1}^k|U_{ir}|^2,\\
 \varepsilon_{ij}&=\frac{W_{ij}}{k}\sum_{r=1}^k
 \left|\frac{U_{ir}}{\sqrt{d_i}}-
 e^{i\theta_{ij}}\frac{U_{jr}}{\sqrt{d_j}}\right|^2.
 \label{eq:energy}
\end{align}
The second expression applies to positive-degree endpoints, counting each
unordered edge once. Mass sums to one; energy sums to the cluster's mean
eigenvalue. Both are invariant to basis choice within the cluster. Effective
support $1/\sum_i m_i^2$ measures concentration, not the number of pages with
nonzero mass. Reproduction sources retain full pairwise results, the synthetic mechanism
and finite-change diagnostics.

\subsection{Panels and Numerical Qualification}
Graph comparisons and numerical checks distinguish structural responses from computation error.
The characterization panel contains one snapshot per corpus. A fixed-text Docker
panel has 32 orientation versions in each of two edge families: as written,
globally reversed, fully reciprocal, twenty random orientations and nine partial
reversals. The sensitivity family removes index and release-note scaffolding.
Except for the reciprocal version, these interventions preserve $W$. They need
not preserve directed degrees. Version-level exposure associations use the 29
random and partial variants, with 20,000 permutations, within-family Holm
correction, and partial correlations controlling sources, sinks, out-degree Gini
and mean seed out-degree.

An exploratory matched study searches the fourteen frozen LWCCs for simple
nonreciprocal directed cycles of lengths 3--6. Seeded hash ordering selects up
to two per length from encountered candidates, with 250,000 search expansions
per length and corpus. Selection precedes reversal spectra; capped searches
are not uniform samples of all cycles. Each reversal starts from its original
baseline, preserving nodes, $W$, directed and weighted degrees, reciprocity,
sources, sinks, Gini and weak/strong partitions. Fixed $W$ and in-minus-out
degrees also preserve trophic incoherence. The study uses $q=.25$ and a fixed $.10$
companion, with $q=0$ equality and global-reversal controls for every corpus.
Separate five-node, path and isolate fixtures test mechanism and blind spots.

The implementation uses complex double precision. Acceptance requires relative
residual and orthogonality errors at most $10^{-10}$, Hermitian error at most
$10^{-12}$, and computed modes within $[0,2]$ to $10^{-10}$. The near-zero
tolerance is $10^{-9}$. Dense and sparse solvers retain raw values; failed
qualification is not an observation. For attribution, a change is resolved when
$|\Delta\lambda_1|>\max(10^{-8},10(r_A+r_B))$, using the absolute ground residuals.
Complete ground clusters and separated boundaries are required.

DocSpectraICSC2027 v1.1.0~\cite{docspectraicsc2027} records the original measurement
and attribution implementations; the extended matched study uses RetrievalArena
revision \texttt{38c8ed5}.
Execution was verified on Linux/amd64; other platforms have not been tested.
